\documentclass[12pt]{article}
\usepackage[utf8]{inputenc}
\usepackage[T1]{fontenc}
\usepackage{hyperref}
\title{Ware Language Resources for K-12 Curriculum Development}
\author{Compiled by Claude Sonnet 5.5 (claude-sonnet-5-5) for Andr\'as Kornai}
\date{}
\begin{document}
\maketitle

\section{Introduction}
The ISO 639-3 code \texttt{wre} (Glottocode ware1252, listed as ``Unattested: Ware'', macro-area Africa) was, according to the sources found, retired effective 2008-01-14 with the stated reason ``non-existent''. The change request reportedly said that survey teams in Tanzania (Mara and Arusha regions, 2005--2006) found no knowledge of a Ware people or language. Wikipedia describes Ware as an extinct Bantu language near Lake Victoria (Guthrie code JE.407). It should not be confused with Wares (\texttt{wai}, ware1254) or Warembori (\texttt{wsa}, ware1253, a language of New Guinea). No evidence was found that a speaker community, or any language material, exists.

\section{Primary Educational Resources}
None found. No primers, textbooks, or school programmes exist in the search results.

\section{Linguistic and Reference Materials}
\begin{itemize}
\item Glottolog entry for Ware (unattested): \url{https://glottolog.org/resource/languoid/id/ware1252}
\item Wikipedia article on the Ware language: \url{https://en.wikipedia.org/wiki/Ware_language}
\end{itemize}
These pages give only classification and status. They contain no word lists or grammar.

\section{Religious and Community Resources}
None found.

\section{Digital and Audio Resources}
None found.

\section{Curriculum Development Context}
The language appears to be non-existent or unattested, so there is no material on which a native-language K-12 curriculum could be built. A curriculum project would first need to establish whether a community exists. Nearby living Bantu languages of the Lake Victoria region might be a more realistic target, but that is outside this survey entry.
\end{document}